\scanpage{146}
\mspage{V.18}
% archive sheet number: 145

sont relativement \deleted{\unsure{fondamentaux}} \inserted{auxiliaires}, contrairement à \deleted{\unsure{celui|celle|ceux}} ceux de $V_s,V,V_{\unsure{\bar\eta|\bar y}}$. \deleted{\unsure{Il semble que}} On peut dire que la \deleted{\unsure{toute la structure|structure topologique}} \deleted{qui contiennent les informations} connaissance de la structure topologique de ces derniers contient \deleted{\unsure{structures topologiques des derniers|structures topologiques des fibres}} l'essentiel de \deleted{\unsure{ce que}} l'étude topologique de $f:X\longrightarrow S$ en le point $x$.

4.3. D'autre part, le groupe $I$ opère sur $\bar V$ \inserted{(3.4), via des opérations \unsure{sur}} par $V$-automorphismes, en induisant des opérations \uline{triviales} sur $V_{\bar s}$, et sur $V_{\bar\eta}$ les \inserted{$V_\eta$-}opérations évidentes sur l'écriture $V_{\bar\eta}\simeq V_\eta\times_\eta\bar\eta$.

4.4. Explicitons maintenant, par des triangles exacts dans des catégories dérivées convenables, les principales relations \deleted{\unsure{\&}} entre les cohomologies des cinq schémas envisagés, en choisissant des coefficients $\Lambda_{\unsure{\nu|v}}=\mathbb Z/\ell^{\nu+1}\mathbb Z$ comme d'habitude. Nous considérons $\mathbb R\unsure{\Gamma|\mathscr T}_{\deleted{\unsure{I|*}}}(\bar V,\Lambda_\nu)$ comme élément de $D(\Lambda_\nu[I])$, catégorie dérivée de la catégorie des $\Lambda_\nu$-modules à action admissible de $I$. Pour rappeler ce fait, nous écrirons plutôt $\mathbb R(\Gamma,I)(\bar V,\Lambda_\nu)$ ; l'application du foncteur oubli à cet objet redonne $\mathbb R\Gamma(\bar V,\Lambda_\nu)$. On a alors un triangle exact

\deleted{\unsure{(4.4.1)}} (A)
\diagram{\begin{tikzcd}[ampersand replacement=\&,column sep=large,row sep=large]
\& \mathbb R(\Gamma,I)(V_{\bar\eta},\Lambda_\nu)\arrow[dl,"{\unsure{\partial|\delta|)}}"']\&\\
\mathbb R(\Gamma,I)(V_{\bar s},\underline\Phi^\bullet_{V_{\bar s}})\arrow[rr] \& \& \mathbb R(\Gamma,I)(\bar V,\Lambda_\nu)\arrow[ul,"\bar j^*"]
\end{tikzcd}},
où $\underline\Phi^\bullet_{V_{\bar s}}=\mathbb R.\Gamma_{V_{\bar s}}(\Lambda_{\nu\bar V})$, peut aussi s'interpréter comme
